\section{Setup and Notation}

% Scope:
% Define the canonical linear-regression/search problem once, before splitting
% into regimes.
%
% Include the data-generating objects:
% treatment or focal regressor x_0, candidate controls X, outcome y, population
% covariance blocks (v_{00}, d, C), signal coefficients beta_0 and beta_{-0},
% and noise scale sigma_epsilon.
%
% Define the search spaces:
% exact-size supports S_s, at-most-size supports S_{\le s}, sparse residual
% makers M_S, and the searched coefficient and t-statistic.
%
% Introduce common quantities used throughout:
% L_{s,p}, sparse overlap eta_s^{(d,C)}, sparse condition number kappa_s(C),
% deterministic sparse drift b_C(S), center Theta_{C,s}, and fluctuation scales.
%
% Keep this section neutral: no regime-specific theorem statements except for
% notation that all later sections need.
